
opérations @@M0@@
définissent d'après 7.3 des endomorphismes additifs du foncteur en @@M1@@-anneaux @@M2@@~; d'après 1.10, ils sont caractérisés par leur valeur pour @@M3@@. On a :
@@M4@@
Par suite :
@@M5@@

\textbf{Proposition 7.5.} i) \uline{Si} @@M6@@ \uline{est un} @@M7@@-\uline{anneau}, \uline{les opérations} @@M8@@ \uline{vérifient} \uline{les identités} :
@@M9@@
ii) \uline{Réciproquement}, \uline{si les} @@M10@@ \uline{vérifient ces identités}, \uline{et si} @@M11@@ \uline{est de} \uline{caractéristique nulle}, @@M12@@ \uline{est un} @@M13@@-\uline{anneau}.

i) Si @@M14@@ est un @@M15@@-anneau, on a @@M16@@, d'où : @@M17@@, et par suite :
@@M18@@

Pour montrer la seconde des relations (7.5.1), on considère le diagramme :
@@M19@@
où les flèches verticales sont définies par la multiplication, et les applications @@M20@@ comme suit : l'application « dérivée logarithmique » est
